% =========================================================================
% Appendix A: Raw triplicate diameters and phytochemical colour keys
% =========================================================================

\chapter{Raw triplicate measurements}
\label{app:data_template}

All diameters below are millimetres and include the disc. They are transcribed from the laboratory summary table (C1 = \qty{50}{\milli\gram\per\milli\litre}; C2 = \qty{100}{\milli\gram\per\milli\litre}).

\section{Escherichia coli}
\label{sec:app_ecoli}

\begin{table}[htbp]
  \centering
  \caption{Raw zones (mm) against \textit{E.\ coli}.}
  \label{tab:app_ecoli}
  \begin{tabular}{lccc}
    \toprule
    Treatment & R1 & R2 & R3 \\
    \midrule
    C1 & 12.00 & 11.78 & 12.05 \\
    C2 & 15.00 & 15.50 & 14.48 \\
    Gentamicin disc & 17.00 & 17.00 & 17.00 \\
    Distilled water & 0.00 & 0.00 & 0.00 \\
    \bottomrule
  \end{tabular}
\end{table}

\section{Pseudomonas aeruginosa}
\label{sec:app_pa}

\begin{table}[htbp]
  \centering
  \caption{Raw zones (mm) against \textit{P.\ aeruginosa}.}
  \label{tab:app_pa}
  \begin{tabular}{lccc}
    \toprule
    Treatment & R1 & R2 & R3 \\
    \midrule
    C1 & 16.00 & 16.30 & 15.80 \\
    C2 & 13.00 & 12.50 & 12.10 \\
    Gentamicin disc & 17.50 & 17.00 & 17.00 \\
    Distilled water & 0.00 & 0.00 & 0.00 \\
    \bottomrule
  \end{tabular}
\end{table}

\section{Klebsiella pneumoniae}
\label{sec:app_kp}

\begin{table}[htbp]
  \centering
  \caption{Raw zones (mm) against \textit{K.\ pneumoniae}.}
  \label{tab:app_kp}
  \begin{tabular}{lccc}
    \toprule
    Treatment & R1 & R2 & R3 \\
    \midrule
    C1 & 12.00 & 13.00 & 12.05 \\
    C2 & 13.50 & 13.50 & 12.75 \\
    Gentamicin disc & 15.50 & 15.00 & 15.00 \\
    Distilled water & 0.00 & 0.00 & 0.00 \\
    \bottomrule
  \end{tabular}
\end{table}

\section{Phytochemical colour key}
\label{sec:app_phyto}

\begin{table}[htbp]
  \centering
  \caption{Colour recorded as positive in the laboratory phytochemical sheet.}
  \label{tab:app_phyto}
  \begin{tabularx}{\textwidth}{
    >{\raggedright\arraybackslash}p{0.16\textwidth}
    >{\raggedright\arraybackslash}p{0.22\textwidth}
    >{\raggedright\arraybackslash}X
    >{\raggedright\arraybackslash}p{0.20\textwidth}
  }
    \toprule
    Constituent & Test & Reagents & Colour / sign \\
    \midrule
    Alkaloids & Mayer's & 2\,mL Mayer's + extract & Dull white precipitate \\
    Carbohydrates & Fehling's & Extract + Fehling A and B, heated & Brick-red precipitate \\
    Saponins & Foam & 1\,mL extract + 5\,mL water, shaken & Foam \\
    Phytosterols & Salkowski & 2\,mL extract + 2\,mL chloroform + 2\,mL \ce{H2SO4} & Reddish brown \\
    Tannins & \ce{FeCl3} & Extract + 3--4 drops \ce{FeCl3} & Dark blue / greenish black \\
    Phenols & \ce{FeCl3} & Extract + 10\% \ce{FeCl3} & Bluish black \\
    Flavonoids & Alkaline reagent & Extract + 10\% \ce{NaOH} & Yellow \\
    Terpenoids & Salkowski (thionyl chloride) & 1\,mL extract + thionyl chloride & Pink \\
    Quinones & Sulphuric acid & 1\,mL extract + 1\,mL \ce{H2SO4} & Red \\
    \bottomrule
  \end{tabularx}
\end{table}

\section{Items not in the laboratory record}
\label{sec:app_missing}

No ATCC numbers, no KASH voucher code, no disc-load volume, no MIC/MBC table, and no \textit{P.\ aeruginosa} inhibition photograph were present in the source files used to compile this appendix.
